\section{A restriction preserving own-question variance}
\label{sec:restriction}

The function from the previous section has variance of order \(d\),
but the functions used to extract labels must depend only on one
position's own questions and weights. We shall freeze all but a batch
of \(h\) positions and retain variance of order \(h\) in conditional
averages with precisely this dependence. Questions at different positions are
correlated because they share pair seeds; our orthogonal decomposition
therefore uses the independent seeds and weights themselves.

Let \(d\geq4\) and \(r,M\geq1\) be integers, and put
\(\mathcal E=\binom{[d]}2\).  For each \(e\in\mathcal E\), draw
\(c_e\) uniformly from \([M]\), independently over \(e\).  At each
endpoint \(j\in e\), this seed reveals a question
\(Q_{e,j}(c_e)\), where \(Q_{e,j}\) is a fixed map to a finite set.
There is no independence assumption on the two questions from the same seed.
In the construction, for \(e=\{j,k\}\) with \(j<k\), these maps are
\(Q_{e,j}(c)=x_c\) and \(Q_{e,k}(c)=y_c\).
Independently of all seeds, draw \(s_j\) uniformly from \([-1,1]^r\)
for each \(j\in[d]\), independently over \(j\).  All expectations and
\(L^2\)-norms in this section initially refer to the resulting product law
on
\[
 \Omega=[M]^{\mathcal E}\times([-1,1]^r)^d.
\]

We first specify the restrictions and conditional averages in the lemma.
Given a bounded Borel measurable function \(f(\mathbf c,\mathbf s)\) and a set
\(J\subseteq[d]\), its \emph{frozen data} are
\begin{equation}\label{res:frozen-data}
 \mathbf a=
 \bigl((c_e)_{e\notin\binom J2},(s_k)_{k\notin J}\bigr).
\end{equation}
Fixing these values leaves a function \(f^*\) of the hidden seeds
\((c_e)_{e\in\binom J2}\) and the weights \((s_j)_{j\in J}\), with
their original product law.  For \(j\in J\), write
\[
 u_j=\bigl(Q_{\{j,k\},j}(c_{\{j,k\}})\bigr)_{k\in J\setminus\{j\}},
\]
with entries in increasing order of \(k\).  Define
\(H_j^{J,\mathbf a}(u,s)\) by inserting \(s_j=s\) in \(f^*\) and
averaging the remaining variables conditional on \(u_j=u\).  Explicitly,
each incident hidden seed \(c_{\{j,k\}}\) is drawn uniformly from the
fiber of \(Q_{\{j,k\},j}\) specified by the corresponding entry of
\(u\); these draws are independent.  The other hidden seeds and other
weights retain their original independent laws.  On a question tuple
of probability zero set \(H_j^{J,\mathbf a}=0\).  This definition uses
ordinary finite averages and cube integrals with the frozen values plugged
in, so it also defines the averages when \(\mathbf a\) has continuous
coordinates.  We abbreviate \(H_j^{J,\mathbf a}\) to \(H_j\) when the
restriction is fixed.

\begin{lemma}[Own-question variance under a restriction]
\label{res:private-variance}
Let \(\nu>0\) and let \(h\) be an integer with \(2\leq h\leq d\).
Suppose the bounded Borel measurable function \(f\) on \(\Omega\) satisfies
\begin{equation}\label{res:variance-hypotheses}
 \E_{\mathbf s}f(\mathbf c,\mathbf s)=0
       \quad\text{for every }\mathbf c,
 \qquad \E f^2\geq\nu d.
\end{equation}
Suppose also that the following oscillation bounds hold with all coordinates
not mentioned kept fixed:
\begin{enumerate}[label=\textup{(\roman*)}]
 \item Changing one weight block \(s_j\) changes \(f\) by at most \(2\).
 \item Changing one seed \(c_e\) changes \(f\) by at most \(4\).
 \item For any \(k\in[d]\), simultaneously changing seeds incident to
 \(k\), while preserving \(Q_{\{k,j\},j}(c_{\{k,j\}})\) for every
 \(j\ne k\), changes \(f\) by at most \(2\).
\end{enumerate}
If
\begin{equation}\label{res:scale}
 16h+32h^2\leq\frac{\nu d}{2},
\end{equation}
then there are a fixed \(h\)-element set \(J\) and fixed frozen data
\(\mathbf a\) such that
\begin{equation}\label{res:variance}
 \sum_{j\in J}\E_{u_j}\Var_{s_j}H_j(u_j,s_j)
       \geq\frac{\nu h}{4}.
\end{equation}
The law in \eqref{res:variance} is the product law of the remaining hidden
seeds and batch weights described above.  In particular it is not conditioned
on an event involving a draw of those remaining variables.
\end{lemma}

\begin{proof}
The conditional mean \(H_j\) averages all unfrozen weights except
\(s_j\). It also averages every hidden seed on a pair avoiding \(j\),
and retains only \(j\)'s question on each hidden pair incident to \(j\).
These are the three ways in which averaging can lose information.
We use the product-space orthogonal decomposition associated with
the Efron--Stein method~\cite[Section~2]{EfronStein1981}, giving the needed
projection identities explicitly. We assign each component to one
weight block on which it depends, and ask when a random batch preserves
that component in its assigned conditional mean. The oscillation
bounds control the energy lost through each type of averaging.

\paragraph{Averaging projections and weight supports.}
Let \(W_j^0\) average the weight block \(s_j\), retaining all other
coordinates, and set \(W_j^1=\operatorname{Id}-W_j^0\).  Let \(P_e^0\)
average the seed \(c_e\).  For \(j\in e\), let \(O_{e,j}\) average
that seed conditional on its question \(Q_{e,j}(c_e)\), retaining all
other coordinates.  Each averaging is an orthogonal projection in
\(L^2(\Omega)\): averaging twice gives the same function, and integrating
over the retained data shows self-adjointness.  Also
\[
 \operatorname{ran}(P_e^0)\subseteq\operatorname{ran}(O_{e,j}),
 \qquad P_e^0O_{e,j}=O_{e,j}P_e^0=P_e^0.
\]
Operators on distinct coordinates commute by the product law and Fubini's
theorem.  The two endpoint projections \(O_{e,j},O_{e,k}\) on a single
seed need not commute; we will not assume that they do.

For \(S\subseteq[d]\), put
\[
 \Pi_S=\prod_{j\in S}W_j^1\prod_{j\notin S}W_j^0,
 \qquad f_S=\Pi_Sf.
\]
Expanding \(\prod_j(W_j^0+W_j^1)\) gives \(f=\sum_Sf_S\).
Distinct \(S\)'s give orthogonal components, because their projections
have a factor \(W_j^0W_j^1=0\) at some position.  The first condition
in \eqref{res:variance-hypotheses} says that \(f_\varnothing=0\).
Furthermore
\(f-W_j^0f=\sum_{S\ni j}f_S\), and hence
\begin{equation}\label{res:weight-degree}
 \sum_{S\ne\varnothing}|S|\|f_S\|_2^2
   =\sum_{j=1}^d\|f-W_j^0f\|_2^2
   \leq4d.
\end{equation}
Indeed, on every fiber obtained by fixing the other coordinates, the
one-weight oscillation bound implies \(|f-W_j^0f|\leq2\).

\paragraph{Refining the seed coordinates.}
For every nonempty \(S\), choose a designated position \(j(S)\in S\)
by a fixed rule, for example \(j(S)=\min S\).  This choice depends only
on the support \(S\).  With \(j=j(S)\), split each seed coordinate into
the following projections:
\[
 \begin{array}{ll}
 j\notin e:
   &P_e^0,\quad\operatorname{Id}-P_e^0,\\[3pt]
 j\in e:
   &P_e^0,\quad O_{e,j}-P_e^0,\quad\operatorname{Id}-O_{e,j}.
 \end{array}
\]
The nested ranges just proved make the three projections in the second
line pairwise orthogonal, and their sum is the identity.  For a choice
\(\omega\) of one displayed projection at every seed, let
\(f_{S,\omega}\) be their product applied to \(f_S\).  For fixed
\(S\), these components give an orthogonal decomposition of \(f_S\).
The seed projections commute with all weight projections, so every
\(f_{S,\omega}\) remains in \(\operatorname{ran}(\Pi_S)\).
Consequently the components for different \(S\)'s are also orthogonal,
even though the endpoint used to refine a seed can change with \(S\).
Thus
\begin{equation}\label{res:component-energy}
 \sum_{S,\omega}\|f_{S,\omega}\|_2^2=\|f\|_2^2,
\end{equation}
where henceforth these sums range over nonempty \(S\) and its allowed
choices \(\omega\).

Let \(T_{S,\omega}\subseteq\mathcal E\) be the seeds where the chosen
projection is not \(P_e^0\).  Let
\(N_{S,\omega}\subseteq[d]\setminus\{j(S)\}\) be the positions
\(k\) for which the projection at \(e=\{j(S),k\}\) is
\(\operatorname{Id}-O_{e,j(S)}\).  These are the incident components
not retained by the designated position's own question.
For each fixed \(e\), applying \(\operatorname{Id}-P_e^0\) keeps
exactly the components with \(e\in T_{S,\omega}\).  Orthogonality and
the one-seed oscillation bound therefore give
\begin{equation}\label{res:seed-degree}
 \sum_{S,\omega}|T_{S,\omega}|\|f_{S,\omega}\|_2^2
  =\sum_{e\in\mathcal E}\|f-P_e^0f\|_2^2
  \leq16\binom d2\leq8d^2.
\end{equation}
Here \(|f-P_e^0f|\leq4\) follows by averaging values whose oscillation
is at most \(4\).

At this stage the ordinary seed-energy bound is of order \(d^2\).
That is sufficient for hidden pairs avoiding the designated position,
because both endpoints must enter the small batch. It is insufficient
for a hidden incident seed, which needs only one extra endpoint in the
batch. The next estimate reduces the relevant incident energy to order
\(d\) by resampling a whole star at once.

\paragraph{Charging the residual components to stars.}
For each \(k\in[d]\), define
\[
 D_k=\prod_{j\ne k}O_{\{k,j\},j}.
\]
The factors act on distinct seed coordinates, so this is an orthogonal
projection.  More explicitly, it retains all weights, all seeds not
incident to \(k\), and the opposite endpoint question at each seed
incident to \(k\).  Given these data, the incident seeds have independent
conditional laws on their respective fibers; their product gives exactly
the displayed operator.  \Cref{fig:star} illustrates the information retained by this operation.
Hypothesis~\textup{(iii)} bounds the oscillation
of \(f\) on each such fiber by \(2\).  It follows that
\begin{equation}\label{res:star-energy}
 \|f-D_kf\|_2^2\leq4.
\end{equation}
\begin{figure}[!htb]
\centering
\begin{tikzpicture}[>=Stealth, every node/.style={font=\small}]
\node[draw,circle,minimum size=9mm,fill=blue!8] (k) at (0,0) {$k$};
\node[draw,rounded corners,align=center] (j) at (-4,1.2) {position $j$\\question $Q_{\{k,j\},j}$ fixed};
\node[draw,rounded corners,align=center] (l) at (4,1.2) {position $j' $\\question $Q_{\{k,j'\},j'}$ fixed};
\node[draw,rounded corners,align=center] (a) at (0,-1.8) {every other position $j''$\\question $Q_{\{k,j''\},j''}$ fixed};
\draw[<->,thick] (k)--node[below,sloped,font=\scriptsize] {$c_{\{k,j\}}$} (j);
\draw[<->,thick] (k)--node[below,sloped,font=\scriptsize] {$c_{\{k,j'\}}$} (l);
\draw[<->,thick] (k)--(a);
\end{tikzpicture}
\caption{The star operation centered at $k$. Each incident seed is
resampled within the fiber of its question at the opposite endpoint.
All nonincident seeds and all weights stay fixed. Thus every query
except the query at $k$ stays fixed. The schematic depicts the
conditional averaging operator $D_k$, not additional graph edges.}
\label{fig:star}
\end{figure}


For a component designated at \(j\), we charge invisible information
on \(\{j,k\}\) to the star centered at \(k\): that star retains
\(j\)'s question on this pair. For every \(j\ne k\), the range
of \(D_k\) is contained in that of
\(O_{\{k,j\},j}\).  Distance to the larger closed subspace is no
greater than distance to the smaller one.  Thus, for every \(v\in L^2\),
\[
 \|(\operatorname{Id}-O_{\{k,j\},j})v\|_2^2
       \leq\|(\operatorname{Id}-D_k)v\|_2^2.
\]
For fixed \(S\) and \(k\ne j(S)\), its seed decomposition shows that
\((\operatorname{Id}-O_{\{k,j(S)\},j(S)})f_S\) is the sum of exactly
the components with \(k\in N_{S,\omega}\).  Consequently
\begin{align}
 \sum_{S,\omega}|N_{S,\omega}|\|f_{S,\omega}\|_2^2
  &=\sum_{S\ne\varnothing}\sum_{k\ne j(S)}
    \|(\operatorname{Id}-O_{\{k,j(S)\},j(S)})f_S\|_2^2
       \notag\\
  &\leq\sum_{k=1}^d\sum_{S\ne\varnothing}
       \|(\operatorname{Id}-D_k)f_S\|_2^2
       \notag\\
  &=\sum_{k=1}^d\|(\operatorname{Id}-D_k)f\|_2^2
    \leq4d.\label{res:residual-degree}
\end{align}
The inequality applies the preceding projection bound and adds the
nonnegative terms with \(j(S)=k\). The next
equality uses the fact that \(D_k\) commutes with the weight projections,
so the images of distinct \(f_S\)'s remain orthogonal.  This argument
uses neither products of different \(D_k\)'s nor commutation of
opposite endpoint projections on one seed.

The three energy budgets are now available: weight interactions cost
at most \(4d\), all nonconstant seeds at most \(8d^2\), and incident
information hidden from the designated question at most \(4d\).
We next check exactly which of these components a restriction preserves.

\paragraph{What the restriction keeps.}
For fixed \(J\) and \(j\in J\), define the projection
\begin{equation}\label{res:restriction-operator}
 K_{J,j}=
 W_j^1\prod_{k\in J\setminus\{j\}}W_k^0
 \prod_{e\in\binom{J\setminus\{j\}}2}P_e^0
 \prod_{k\in J\setminus\{j\}}O_{\{j,k\},j}.
\end{equation}
All seed factors here act on distinct coordinates.  Before applying
\(W_j^1\), these operators average exactly the unfrozen data other than
the own questions \(u_j\) and the own weights \(s_j\).  They retain
the frozen data \(\mathbf a\) as variables.  The final centering
\(W_j^1\) subtracts the mean in \(s_j\), which is independent of
the retained seed information.  Thus
\[
 K_{J,j}f
  =H_j^{J,\mathbf a}(u_j,s_j)
    -\E_{s'_j}H_j^{J,\mathbf a}(u_j,s'_j)
\]
as functions on the original product space.  If
\[
 \mathcal V(J,\mathbf a)=
     \sum_{j\in J}\E_{u_j}\Var_{s_j}H_j^{J,\mathbf a}(u_j,s_j),
\]
then integration over the initially random frozen data gives
\begin{equation}\label{res:restriction-variance}
 \E_{\mathbf a}\mathcal V(J,\mathbf a)
     =\sum_{j\in J}\|K_{J,j}f\|_2^2.
\end{equation}
This identity uses only independent seed coordinates.  Conditioning on
the own questions restricts each incident hidden seed to its own fiber,
as in the definition of \(H_j\).

Consider \(f_{S,\omega}\) and its designated position \(j=j(S)\).
This component is preserved by the term \(K_{J,j}\) whenever
\begin{equation}\label{res:survival}
 S\cap J=\{j\},\qquad
 N_{S,\omega}\cap J=\varnothing,\qquad
 T_{S,\omega}\cap\binom{J\setminus\{j\}}2=\varnothing.
\end{equation}
To verify that its squared norm can be counted in
\eqref{res:restriction-variance}, first note that \(K_{J,j}\) commutes
with every \(\Pi_S\).  Thus its outputs from different weight supports
are orthogonal, and
\(\|K_{J,j}f\|_2^2=\sum_S\|K_{J,j}f_S\|_2^2\).
The weight factors in \(K_{J,j}\) keep \(f_S\) precisely when
\(S\cap J=\{j\}\).  For those \(S\) whose designated position is
this \(j\), the seed factors are diagonal in the chosen seed
decomposition: \(P_e^0\) on a nonincident hidden seed keeps just its
constant component, and \(O_{e,j}\) on an incident hidden seed keeps
its constant and own-question components.  The other seed coordinates
are untouched.  They therefore keep exactly the components satisfying
the last two conditions of \eqref{res:survival}, with no cancellation
among them.  We may discard all other nonnegative squared norms in
\eqref{res:restriction-variance} and count each component solely in its
designated term.  It follows that
\begin{equation}\label{res:surviving-energy}
 \E_{\mathbf a}\mathcal V(J,\mathbf a)
   \geq\sum_{S,\omega}
       \mathbf1_{\{\text{conditions \eqref{res:survival} hold}\}}
       \|f_{S,\omega}\|_2^2.
\end{equation}

\paragraph{Averaging the restriction.}
Choose \(J\) uniformly among all \(h\)-element subsets of \([d]\).
For a fixed component, the designated position belongs to \(J\) with
probability \(h/d\).  Conditional on this inclusion, a specified other
position belongs to \(J\) with probability
\[
 p_1=\frac{h-1}{d-1},
\]
and a specified pair avoiding the designated position lies in \(J\)
with probability
\[
 p_2=\frac{(h-1)(h-2)}{(d-1)(d-2)}.
\]
The first two conditions in \eqref{res:survival} fail only if one of
\(S\setminus\{j(S)\}\) or \(N_{S,\omega}\) is included.  The third
can fail only at a pair in \(T_{S,\omega}\) avoiding \(j(S)\).
A union bound therefore gives conditional survival probability at least
\begin{equation}\label{res:survival-probability}
 1-p_1\bigl(|S|-1+|N_{S,\omega}|\bigr)-p_2|T_{S,\omega}|.
\end{equation}
Charging all of \(T_{S,\omega}\) only decreases this lower bound.
The bound remains valid when the displayed expression is negative.

By \eqref{res:weight-degree}, \eqref{res:residual-degree}, and the
orthogonal refinement of each \(f_S\),
\[
 \sum_{S,\omega}\bigl(|S|-1+|N_{S,\omega}|\bigr)
             \|f_{S,\omega}\|_2^2\leq8d.
\]
Combining \eqref{res:component-energy}, \eqref{res:seed-degree},
\eqref{res:surviving-energy}, and \eqref{res:survival-probability}
now yields
\begin{align}
 \E_{J,\mathbf a}\mathcal V(J,\mathbf a)
   &\geq\frac hd\bigl(\|f\|_2^2-8dp_1-8d^2p_2\bigr)\notag\\
   &\geq\frac hd\bigl(\nu d-16h-32h^2\bigr)
    \geq\frac{\nu h}{2}.\label{res:averaged-variance}
\end{align}
For the second inequality, \(d\geq4\) gives
\(p_1\leq2h/d\) and \(p_2\leq4h^2/d^2\); the last inequality
is \eqref{res:scale}.

The nonnegative measurable quantity \(\mathcal V(J,\mathbf a)\) is
bounded, since \(f\) is bounded.  Its mean in
\eqref{res:averaged-variance} implies that a set of choices
\((J,\mathbf a)\) of positive probability has
\(\mathcal V(J,\mathbf a)\geq\nu h/4\).  Choose one such restriction.
The conditional averages were defined by explicit integration with its
frozen coordinates plugged in.  All hidden seeds and batch weights are
then drawn with their original product law, establishing
\eqref{res:variance} with the asserted interpretation.
\end{proof}
